\chapter{Diseño del software}

\section{Conceptos y principios de diseño}

El diseño es la actividad que decide cómo se organiza el sistema para satisfacer
los requisitos. Entre la especificación y el código hay un espacio de decisiones
que no está determinado por ninguno de los dos, y es ahí donde se juega la
mantenibilidad futura.

Los conceptos que sostienen esas decisiones son pocos y llevan décadas siendo los
mismos \cite{pressman2021}:

\begin{itemize}
    \item \textbf{Abstracción.} Trabajar con una descripción que retiene lo
          relevante y omite el resto, y hacerlo por niveles.
    \item \textbf{Modularidad.} Dividir el sistema en partes con una
          responsabilidad identificable, de forma que cada una pueda entenderse y
          modificarse por separado.
    \item \textbf{Ocultación de información.} Que cada módulo esconda sus
          decisiones internas tras una interfaz. Es lo que permite cambiar la
          implementación sin arrastrar a los clientes del módulo.
    \item \textbf{Separación de intereses.} Que cada parte se ocupe de una
          cuestión y no de varias entremezcladas.
\end{itemize}

De ellos se derivan las dos medidas con las que se juzga una descomposición:

\begin{itemize}
    \item La \textbf{cohesión} mide hasta qué punto los elementos de un módulo
          contribuyen a un mismo propósito. Interesa que sea alta.
    \item El \textbf{acoplamiento} mide la dependencia entre módulos. Interesa que
          sea bajo, y sobre todo que las dependencias vayan hacia abstracciones
          estables y no hacia detalles.
\end{itemize}

Las dos van juntas: una descomposición que reparte una misma responsabilidad entre
varios módulos baja la cohesión y sube el acoplamiento a la vez, porque esos
módulos tendrán que cambiar siempre juntos.

Por encima de estas medidas están los principios que guían la asignación de
responsabilidades. Los patrones GRASP —experto en información, creador,
controlador, bajo acoplamiento, alta cohesión— dan un vocabulario para justificar
por qué una operación va en una clase y no en otra, que es la decisión más
frecuente del diseño orientado a objetos \cite{larman2006}.

\section{Diseño de los casos de uso}

El diseño de casos de uso conecta el análisis con la estructura de objetos:
partiendo de las operaciones del sistema identificadas en el análisis, se decide
qué objetos colaboran para realizarlas y con qué mensajes.

El instrumento habitual es el \textbf{diagrama de interacción}, en cualquiera de
sus dos formas. El diagrama de secuencia ordena los mensajes en el tiempo y se lee
bien cuando el flujo importa; el diagrama de comunicación destaca los enlaces
entre objetos y se lee mejor cuando lo que interesa es la estructura de
colaboración \cite{arlow2006}.

Diseñar una interacción obliga a repartir responsabilidades, y ese reparto es el
que determina el acoplamiento del resultado. Dos criterios ayudan a decidirlo:

\begin{itemize}
    \item Asignar cada responsabilidad a la clase que ya tiene la información
          necesaria para cumplirla, en vez de crear una clase que pida esa
          información a todas las demás.
    \item Introducir un \emph{controlador} que reciba los eventos del sistema y
          delegue, en lugar de dejar que la capa de presentación hable
          directamente con el dominio.
\end{itemize}

Cuando un objeto necesita datos de muchos otros para hacer su trabajo, la
interacción lo delata: el diagrama se llena de mensajes de consulta salientes.
Esa forma es la señal de que la responsabilidad está mal colocada.

\section{Diseño de la estructura de objetos}

El resultado del diseño de las interacciones se consolida en el \textbf{diagrama
de clases de diseño}, que añade al modelo conceptual lo que el análisis
deliberadamente omitía: las operaciones, la visibilidad de los atributos, los
tipos, la navegabilidad de las asociaciones y las clases que no pertenecen al
dominio sino a la solución.

Hay que distinguirlo del modelo conceptual del que procede. Aquel describía el
problema; este describe la solución, y por eso puede contener clases que no
existen en el dominio —repositorios, adaptadores, fábricas— sin que eso sea un
defecto.

Sobre esa estructura se aplican los \textbf{patrones de diseño}, soluciones
probadas a problemas recurrentes. Los que más se emplean en este nivel son los de
creación, para desacoplar al cliente de la clase concreta que instancia; los
estructurales, como adaptador o fachada, para acomodar interfaces que no encajan;
y los de comportamiento, como estrategia u observador, para dejar variable lo que
se sabe que va a cambiar \cite{larman2006}.

Los patrones se aplican cuando el problema que resuelven está presente. Aplicarlos
por anticipado, «por si acaso», añade indirección sin beneficio y es una forma
frecuente de complicar un diseño sin necesidad.

\section{Arquitectura del software}

La arquitectura es el conjunto de decisiones de diseño de mayor alcance: la
descomposición del sistema en partes principales, la relación entre ellas y las
propiedades globales que esa organización garantiza. Son las decisiones más caras
de revertir, porque el resto del diseño se apoya en ellas.

La arquitectura la determinan, sobre todo, los requisitos no funcionales. El
rendimiento empuja hacia componentes gruesos y poca comunicación entre ellos; la
seguridad, hacia capas con control de acceso en las fronteras; la disponibilidad,
hacia la redundancia; la mantenibilidad, hacia componentes pequeños y
reemplazables. Como estas fuerzas tiran en direcciones opuestas, no existe una
arquitectura buena en abstracto, solo una adecuada a unas prioridades declaradas
\cite{sommerville2016}.

Los estilos más comunes son:

\begin{itemize}
    \item \textbf{En capas.} Cada capa usa solo servicios de la inmediatamente
          inferior. Facilita sustituir una capa entera y aísla el dominio de la
          tecnología, a costa de rendimiento en las llamadas que la atraviesan.
    \item \textbf{Cliente-servidor.} Separa a quien pide de quien sirve, lo que
          permite escalar cada lado por separado.
    \item \textbf{Modelo-vista-controlador.} Aísla la presentación del estado y de
          la lógica, de modo que varias vistas puedan compartir un mismo modelo.
    \item \textbf{Por componentes o servicios.} Divide el sistema en unidades
          desplegables de forma independiente. Gana en escalabilidad y en
          autonomía de los equipos; paga con complejidad de operación y con los
          problemas propios de la comunicación en red.
    \item \textbf{Tuberías y filtros.} Encadena transformaciones sobre un flujo de
          datos. Encaja bien en procesamiento por lotes.
\end{itemize}

Decidir la arquitectura incluye documentarla. Una arquitectura que solo existe en
la cabeza de quien la diseñó deja de cumplirse en cuanto el equipo crece, y la
erosión resultante es difícil de revertir porque el código ya depende de la forma
degradada.
